\section{Comparative Analysis: Claude Code, OpenClaw, and Hermes Agent}
\label{sec:compare}

The preceding sections documented Claude Code's answers to recurring design questions about loop architecture, safety, extensibility, context management, delegation, and persistence. To situate these findings within the broader agent design space, this section compares Claude Code with two independent open-source AI agent systems that answer many of the same design questions from fundamentally different starting points. OpenClaw is a local-first WebSocket gateway that connects messaging surfaces (WhatsApp, Telegram, Slack, Discord, Signal, and others) to an embedded agent runtime, with companion apps on macOS, iOS, and Android~\citep{openclaw2026}. Hermes Agent is a single Python process whose role is set by the entry point under which it was invoked, with one persistence layer and many surfaces fronting one runtime~\citep{hermesagent2026}.\footnote{Source-grounded reading of Hermes surfaced two cases where the project's own documentation has drifted from the implementation: \file{SECURITY.md} documents an approval-mode taxonomy of \texttt{on}/\texttt{auto}/\texttt{off} while \file{tools/approval.py}:721-750 returns \texttt{manual}/\texttt{smart}/\texttt{off}, and both \file{SECURITY.md} and \file{AGENTS.md} describe a delegation depth default of two while \file{tools/delegate\_tool.py}:128 sets \texttt{MAX\_DEPTH = 1}. We cite the source rather than the docs.} The Hermes distribution exposes three console scripts: \texttt{hermes} for the CLI and gateway controller, \texttt{hermes-agent} for batch runs, and \texttt{hermes-acp} for the Agent Client Protocol adapter that lets external IDEs host Hermes the way Hermes's own gateway can host other agents. Where Claude Code is a CLI coding harness bound to a single repository session, OpenClaw is a persistent control plane for multi-channel personal assistance, and Hermes is one process whose role and surface set are determined by which entry point launched it. The three systems occupy different regions of the agent design space. The value of the comparison lies in showing how the same recurring questions produce different architectural answers when the deployment context changes.

\subsection{Six Comparison Dimensions}
\label{sec:compare:dimensions}

\Cref{tab:comparison-dimensions} summarizes the comparison across six dimensions. Each dimension corresponds to a design question that all three systems must answer. Hermes capabilities that fall outside these six dimensions are discussed separately in \Cref{sec:compare:hermes}.

\begin{table*}[h!]
\centering
\caption{Architectural comparison: Claude Code vs.\ OpenClaw vs.\ Hermes Agent across six design dimensions. Each row captures a recurring design question and the different answers the three systems provide.}
\label{tab:comparison-dimensions}
\scriptsize
\setlength{\tabcolsep}{4pt}
\begin{tabularx}{\textwidth}{c>{\raggedright\arraybackslash}X>{\raggedright\arraybackslash}X>{\raggedright\arraybackslash}X}
\toprule
\textbf{Dimension} & \textbf{Claude Code} & \textbf{OpenClaw} & \textbf{Hermes Agent} \\
\midrule
System scope & CLI/IDE coding harness, ephemeral per-session process & Persistent WS gateway daemon, multi-channel control plane & Single Python process, role set by entry point (\texttt{hermes}/\texttt{hermes-agent}/\texttt{hermes-acp}); in-tree plus plugin-tier messaging surfaces; two SQLite files (\file{state.db}, \file{kanban.db}) for persistence \\
\midrule
Trust model & Deny-first per-action rule evaluation with hooks and optional ML classifier; 7 permission modes; graduated trust spectrum & Single trusted operator per gateway; DM pairing and allowlists for inbound channels; opt-in sandboxing with configurable scope (per-agent, per-session, or shared) and multiple backends & Per-action approval in one file (\file{tools/approval.py}); 3 modes (\texttt{manual}, \texttt{smart}, \texttt{off}) plus unconditional \texttt{HARDLINE\_PATTERNS} floor; same approval rendered across CLI, gateway keyboards (Telegram, Discord, QQBot), ACP \texttt{request\_permission}, and a cron-only unattended mode \\
\midrule
Agent runtime & Iterative async generator (\func{queryLoop}) as system center & Embedded \texttt{pi-coding-agent} SDK runner inside gateway RPC dispatch; per-session queue serialization (with optional global lane) & Synchronous \texttt{while} loop inside \texttt{AIAgent.run\_conversation} (\file{run\_agent.py}); iteration count default 90 plus an iteration budget, with a tool-stripped summary call on exhaustion; tool files self-register via \texttt{registry.register()}; parallel tool calls with sequential fallback \\
\midrule
Extension architecture & 4 mechanisms at graduated context costs: MCP, plugins, skills, hooks & Manifest-first plugin system with 12 capability types and central registry; separate skills layer; built-in MCP via \texttt{openclaw mcp} (server and outbound client registry) & 5 surfaces: 3 at Claude-Code level (\file{plugins/<name>/} with lifecycle hooks, bundled \file{skills/}, MCP servers in \file{config.yaml}) plus 2 backend-swap surfaces (\file{plugins/memory/<name>/}, \file{plugins/model-providers/<name>/}); hooks are Python callbacks or external shell commands \\
\midrule
Memory and context & CLAUDE.md 4-level hierarchy; 5-layer compaction pipeline; LLM-based memory scan & workspace bootstrap files (AGENTS.md, SOUL.md, TOOLS.md, IDENTITY.md, USER.md, plus conditionally BOOTSTRAP.md, HEARTBEAT.md, and MEMORY.md); separate memory system (MEMORY.md, daily notes, optional DREAMS.md); auto-compaction with pluggable providers; optional hybrid search (vector + keyword, conditional on embedding provider); experimental dreaming for long-term promotion & One auxiliary-LLM summarizer with token-budget tail protection and a tool-output prune pass; summary prefixed with a ``reference only'' preamble; pre-injection scan against context files (\file{AGENTS.md}, \file{.cursorrules}, \file{SOUL.md}) for injection patterns and invisible Unicode; persistent \file{MEMORY.md}, \file{USER.md} remain authoritative across compaction \\
\midrule
Multi-agent and routing & Task-delegating subagents (e.g., Explore, Plan, general-purpose); worktree isolation; final response text returned to parent & Two separate concerns: (a) multi-agent routing with isolated agents, distinct workspaces, and binding-based channel dispatch; (b) sub-agent delegation with configurable nesting depth (max 5, default 1, recommended 2) and thread-bound sessions & \texttt{delegate\_task} spawns child \texttt{AIAgent} in a thread pool; concurrency cap 3 children default, depth cap 1 (clamped to [1,3]); leaf children cannot delegate by default; children run with persistent memory disabled; Kanban subsystem (SQLite-backed work queue) with stale-claim reclamation for multi-process worker coordination \\
\bottomrule
\end{tabularx}
\end{table*}

\paragraph{System scope and deployment model.}
Claude Code runs as an ephemeral CLI process bound to a single repository. Each session starts and ends with the terminal. OpenClaw runs as a persistent daemon (default port 18789, loopback-only) that owns all messaging surface connections and coordinates clients, tools, and device nodes over a typed WebSocket protocol. Hermes Agent runs as a single long-lived Python process whose role is fixed by the entry point used to launch it; its gateway sub-command drives messaging surfaces through a two-tier adapter system (in-tree adapters under \file{gateway/platforms/} plus a plugin tier under \file{plugins/platforms/} for Google Chat, IRC, and Teams), and it persists session and multi-agent-board state to two on-disk SQLite files (\file{state.db} and \file{kanban.db}) rather than to a typed WebSocket protocol. This places Hermes closer to OpenClaw than to Claude Code in trust topology but closer to Claude Code than to OpenClaw in process topology; the design question of where to put the broker is answered, in Hermes, with no broker at all. These differences in system scope are the most fundamental architectural divergences in the comparison: they determine how every other design question is framed. A compositional relationship also exists across all three: OpenClaw can host Claude Code, OpenAI Codex, and Gemini CLI as external coding harnesses through its ACP (Agent Client Protocol) integration, and Hermes Agent sits on both sides of the ACP host/guest split, with its gateway hosting platform adapters and tool servers while its \texttt{hermes-acp} adapter lets external IDEs host Hermes, making the systems stackable rather than purely alternative.

\paragraph{Trust model and security architecture.}
The three systems address different threat models. Claude Code assumes an untrusted model operating within a trusted developer's machine: the deny-first permission system (\Cref{sec:auth}) evaluates every tool invocation, the ML classifier provides automated safety assessment, and seven permission modes create a graduated autonomy spectrum. OpenClaw assumes a single trusted operator per gateway instance. Its security architecture begins with identity and access control (DM pairing codes, sender allowlists, gateway authentication) rather than per-action safety classification. Tool policy uses configurable allow/deny lists per agent rather than a centralized classifier. Sandboxing is available as an opt-in feature with multiple backends (Docker, SSH, or OpenShell) and configurable scope (per-agent, per-session, or shared); a \texttt{non-main} mode can sandbox all non-main sessions when enabled, though sandboxing is not active by default. The OpenClaw security documentation explicitly states that hostile multi-tenant isolation on a shared gateway is not a supported security boundary. Hermes Agent takes a third position. Tool authorization is centralized in one file (\file{tools/approval.py}), exposing three modes (\texttt{manual}, \texttt{smart}, \texttt{off}) and an unconditional \texttt{HARDLINE\_PATTERNS} floor that blocks \texttt{rm -rf /}, \texttt{mkfs}, raw \texttt{dd} to block devices, fork bombs, and system shutdown commands regardless of mode~\tierB{}. The smart mode runs an auxiliary-LLM risk assessment; an in-source comment names OpenAI Codex's smart approvals as the design inspiration. Because Hermes runs across many surfaces, the same approval flow is rendered four times: a CLI prompt, gateway approval keyboards on Telegram, Discord, and QQBot, an ACP \texttt{request\_permission} round-trip for IDE clients, and a cron-only mode for unattended runs. Compared with Claude Code's seven-mode taxonomy and ML-classifier-driven \texttt{auto} mode, Hermes has a smaller mode taxonomy but the same multi-layer structure: unconditional safety floor, default interactive prompt, opt-in auxiliary-LLM ``smart'' path. The architectural answer is the same; the surface count differs. The difference across the three reflects where the trust boundary sits: Claude Code places it between the model and the execution environment, OpenClaw places it at the gateway perimeter, and Hermes sits between the two with per-action approvals like Claude Code but rendered through many surfaces like OpenClaw.

\paragraph{Agent runtime and tool orchestration.}
All three systems implement agentic loops, but the loops occupy different positions in their respective architectures. In Claude Code, the \func{queryLoop} async generator (\Cref{sec:turn}) is the system's center: all interfaces feed into it, and it directly manages context assembly, model calls, tool dispatch, and recovery. In OpenClaw, the agent runtime is an embedded \texttt{pi-coding-agent} core (an external SDK from the \texttt{pi-mono} project that OpenClaw integrates), sitting inside a larger gateway dispatch layer. The gateway's \texttt{agent} RPC validates parameters, resolves sessions, and returns immediately; the embedded runner then executes the agentic loop while emitting lifecycle and stream events back through the gateway protocol. Runs are serialized through per-session queues and an optional global lane, preventing tool and session races across the multi-channel surface. In Hermes Agent, the loop is a synchronous \texttt{while} inside \texttt{AIAgent.run\_conversation} (\file{run\_agent.py}), gated on an iteration count (default 90) and a separate iteration budget; when the budget is exhausted, a single tool-stripped summary call lets the agent deliver a closing message rather than terminate mid-thought~\tierB{}; each iteration sends tool schemas collected by \texttt{model\_tools.get\_tool\_definitions()} and routes any tool calls through \texttt{model\_tools.handle\_function\_call}, with tool files in \texttt{tools/*.py} self-registering at import time so that adding a tool requires only one \texttt{registry.register()} call, and a single assistant message returning multiple tool calls is executed in parallel with a sequential fallback. All three systems follow the ReAct pattern~\citep{yao2022react}, with OpenClaw's loop a component within a control plane rather than the control plane itself, and Hermes's loop a peer to Claude Code's at the runtime level but shape-different in implementation, a synchronous Python loop with thread-managed concurrency rather than an async generator that yields.

\paragraph{Extension architecture.}
Claude Code's four extension mechanisms (MCP, plugins, skills, hooks) are organized by context cost (\Cref{sec:ext}): hooks consume zero context, skills consume low context, and MCP servers consume high context. All four extend a single agent's context window and tool surface. OpenClaw uses a manifest-first plugin system with four architectural layers (discovery, enablement, runtime loading, surface consumption) and twelve capability types including text inference, speech, media understanding, image/music/video generation, web search, and messaging channels. Plugins register capabilities into a central registry; the gateway reads the registry to expose tools, channels, provider setup, hooks, HTTP routes, CLI commands, and services. OpenClaw also has a separate skills layer with multiple sources (workspace, project-level, personal, managed, bundled, and extra directories, with workspace skills taking highest precedence) plus a public registry (ClawHub) and supports MCP through built-in \texttt{openclaw mcp} commands (server and outbound client registry). Hermes Agent exposes five extension surfaces against Claude Code's four~\tierB{}: three sit at the same level as Claude Code's (general plugins under \file{plugins/<name>/} with lifecycle hooks, bundled skills under \file{skills/}, and MCP servers configured in \file{config.yaml}), while two extra surfaces extend a different axis (pluggable memory providers under \file{plugins/memory/<name>/} and pluggable model providers under \file{plugins/model-providers/<name>/}) that let Hermes swap backends rather than intercept events. Hooks in Hermes can be in-process Python callbacks or config-driven external shell commands, both dispatched through the same hook manager. When Hermes spawns a stdio MCP server, the child's environment is restricted to a small allowlist of system variables before launch, and \texttt{npx}/\texttt{uvx} package names are screened against the OSV malware database. The architectural difference across the three is where the extension acts: Claude Code's extensions modify one agent's action surface, OpenClaw's plugins extend the gateway's capability surface across all agents, and Hermes's extensions add a second axis on top of a Claude-Code-style per-agent surface, swapping entire memory and model-provider backends in place rather than intercepting events around them.

\paragraph{Memory, context, and knowledge management.}
All three systems use transparent file-based memory rather than opaque databases. Claude Code loads a four-level CLAUDE.md hierarchy and manages context pressure through a five-layer compaction pipeline (\Cref{sec:context}). Memory retrieval uses an LLM-based scan of file headers. OpenClaw injects workspace bootstrap files into the system prompt at session start: five core files (AGENTS.md, SOUL.md, TOOLS.md, IDENTITY.md, USER.md) plus conditionally BOOTSTRAP.md, HEARTBEAT.md, and MEMORY.md, with large files truncated. Separately, the memory system manages three file types: \texttt{MEMORY.md} for long-term durable facts, date-stamped daily notes (\texttt{memory/YYYY-MM-DD.md}), and an optional \texttt{DREAMS.md} for dreaming sweep summaries. When an embedding provider is configured, memory search uses hybrid retrieval combining vector similarity with keyword matching. An experimental dreaming system performs background consolidation, scoring candidates and promoting only qualified items from short-term recall into long-term memory. Before compaction, OpenClaw automatically reminds the agent to save important notes to memory files, preventing context loss. Compaction in Hermes is one auxiliary-LLM summarizer with token-budget tail protection, preceded by a tool-output prune pass; the summary is prefixed with a long ``reference only'' preamble that tells the agent the compaction is background context, not new instructions, and that persistent memory (\file{MEMORY.md}, \file{USER.md}) remains authoritative~\tierB{}. Where Claude Code's CLAUDE.md hierarchy is fixed and its compaction has five layers, Hermes has fewer compaction layers but adds a context-file pre-injection scan: files such as \file{AGENTS.md}, \file{.cursorrules}, and \file{SOUL.md} are checked against a list of injection patterns and invisible Unicode characters, and any hit replaces the file's content with a \texttt{[BLOCKED: ...]} marker. All three share the design commitment to user-visible, editable memory but give different answers to the same context-management question: Claude Code invests in graduated context compression (five layers with cache awareness); OpenClaw invests in structured long-term memory promotion (dreaming, daily notes, memory search) plus pluggable compaction providers and session pruning, though its compaction pipeline is less graduated than Claude Code's five-layer system; Hermes invests in pre-injection content scanning and prompt-level isolation of persistent memory from compacted history rather than a deeper compaction stack.

\paragraph{Multi-agent architecture and routing.}
This dimension reveals the starkest architectural difference. Claude Code's multi-agent model is task delegation: the parent spawns subagents (Explore, Plan, general-purpose, and custom types) that operate in isolated context windows with restricted tool sets and return summary-only results (\Cref{sec:subagent}). Worktree isolation provides filesystem-level separation. OpenClaw separates two distinct concerns. First, multi-agent routing: a single gateway can host multiple fully isolated agents, each with its own workspace, authentication profiles, session store, and model configuration, routed to specific channels or senders via deterministic binding rules. Second, sub-agent delegation: within a single agent, background runs can be spawned with configurable nesting depth (maximum 5, default 1, recommended 2), thread-bound sessions on supported channels, and configurable tool policy by depth. OpenClaw's project vision explicitly rejects agent-hierarchy frameworks as a default architecture. Hermes Agent splits the question along yet a different seam. Per-turn delegation goes through the \texttt{delegate\_task} tool, which spawns child \texttt{AIAgent} instances in a thread pool; the parent blocks until children return summaries, concurrency is capped at three children by default and depth is capped at one (clamped to the range [1, 3]), leaf children cannot themselves delegate by default (nested delegation is opt-in through orchestrator children and a higher spawn-depth setting), and children always run with persistent memory disabled so they cannot read or write the parent's notes~\tierB{}. Beyond per-turn delegation, v0.13's Kanban subsystem is a SQLite-backed work queue where multiple Hermes worker profiles claim, heartbeat, and complete tasks, with stale-claim reclamation and an automatic block after two consecutive failed attempts (\file{hermes\_cli/kanban\_db.py}:2703). The distinction across the three matters because Claude Code's subagents are subordinate workers within one user's coding session, OpenClaw's multi-agent routing creates genuinely independent agent instances serving different users or purposes through different channels, and Hermes combines a strict per-turn delegation default with a separate Kanban layer that lets multiple agent processes coordinate over a durable shared board, a question Claude Code and OpenClaw do not ask in the same form.

\subsection{Hermes Agent: Capabilities Beyond the Comparison Table}
\label{sec:compare:hermes}

Session persistence in Hermes uses one WAL-mode SQLite database with full-text search across all session messages and a parent-session chain that records compaction-triggered splits~\tierB{}. Where Claude Code persists sessions to JSONL files per project (\Cref{sec:persist}), Hermes's choice provides cross-session search and concurrent-reader support out of the box. When the gateway restarts mid-conversation, sessions auto-resume on next message arrival, driven by transcript-timestamp freshness logic with a one-hour default window. A built-in cron scheduler and webhook subscription system let the agent run unattended on schedule expressions or external HTTP events; cron sessions run with persistent memory disabled and a configurable inactivity timeout that defaults to 600 seconds. The distribution also includes a background ``curator'' loop that reviews agent-created skills, archives stale ones, and never deletes, gated to skills with agent provenance so bundled and hub-installed skills are off-limits. None of these capabilities map cleanly onto the six comparison dimensions used in the main table, which is why they appear here as a Hermes-specific addendum rather than as cells in \Cref{tab:comparison-dimensions}. OpenClaw's distinctive capabilities (the ClawHub plugin registry under extension architecture and the dreaming memory subsystem under memory and context) already occupy cells in those six dimensions and so do not require a parallel addendum. The Hermes additions in this subsection follow from its persistent-process scope, a deployment-context property Hermes shares with OpenClaw but that Claude Code's per-session CLI model does not provide.

\subsection{What the Contrast Reveals}
\label{sec:compare:reveals}

The comparison surfaces three observations about the design space of AI agent systems.

First, the recurring design questions identified in \Cref{sec:arch:principles} (where reasoning lives, what safety posture to adopt, how to manage context, how to structure extensibility) apply beyond coding agents. OpenClaw answers every one of these questions from the starting point of a multi-channel personal assistant rather than a repository-bound coding tool, and Hermes Agent answers them from a single-process, multi-surface deployment rather than a CLI harness or a multi-channel gateway. The questions are stable; the answers vary with deployment context across the three systems.

Second, the systems make different bets along several dimensions. Claude Code invests in graduated per-action safety evaluation; OpenClaw invests in perimeter-level identity and access control; Hermes sits between the two, with per-action approvals like Claude Code rendered through many surfaces like OpenClaw. Claude Code treats the agent loop as the architectural center; OpenClaw treats the gateway control plane as the center and embeds the agent loop as one component; Hermes's loop is a peer to Claude Code's at the runtime level, occupying the same architectural position rather than sitting inside a gateway dispatch layer. Claude Code's extensions modify a single context window; OpenClaw's plugins extend a shared gateway surface; Hermes adds a second axis on top of a Claude-Code-style per-agent surface, swapping memory and model backends in place rather than intercepting events. These differences are not arbitrary: they follow from the three systems' different trust models and deployment topologies.

Third, the compositional relationships across the three systems are architecturally significant. OpenClaw can host Claude Code as an external coding harness via ACP, and Hermes Agent sits on both sides of the host/guest split: its gateway hosts platform adapters and tool servers, and its \texttt{hermes-acp} adapter lets external IDEs host Hermes. The systems are composable rather than exclusive alternatives, which suggests that the design space of AI agents is not a flat taxonomy but a layered one, where gateway-level systems and task-level harnesses can compose at multiple positions.
